\documentclass[a4paper]{article}
%To  colorize table
\usepackage[table]{xcolor}
\usepackage[a4paper, total={7in, 11.0in},includefoot, includehead, headsep=24pt, headheight=4cm]{geometry}
\usepackage{layout}%No need for this package at production.
\usepackage{setspace}%For switch between line spacing.
\usepackage{graphicx}
%For drawing lines and arbitrary shapes:
\usepackage{tikz}
%To set fixed length for table: 
\usepackage{array}
\usepackage{ifthen}
%To draw squares:
\usepackage{amssymb}
%To trim strings:
\usepackage{trimspaces}
\usepackage{ifthen}
%To draw dynamic multi-column table:
\usepackage{etoolbox}
%\usepackage{biblatex}
%\addbibresource{refs.bib} % with extension
%\setcitestyle{authoryear,open={(},close={)}}
%\bibliographystyle{ieeetr-fa.bst}
%For references:
%\usepackage{biblatex}

%\renewcommand{\bibsection}{}{\rl{{مراجع}\hfill}} 


%\addbibresource{references.bib}


%List of needed counters:
\newcounter{length}
\newcounter{itemCount}


\title{Proposal}
\author{Name}
\date{\today}

\newcommand{\blankField}{ $\ldots\ldots\ldots\ldots\ldots$}
\makeatletter
\newcommand{\trim}[1]{\trim@spaces@noexp{#1}}
\makeatother


\usepackage{fancyhdr}
\pagestyle{fancy}
\fancyhf{}
\usepackage{xepersian}
\settextfont[Scale=1.2]{Yas}
\setlatintextfont[Scale=1]{Times New Roman}
%\bibliographystyle{ieeetr-fa.bst}
%\bibliographystyle{your-bst-file} % Replace 'your-bst-file' with the actual BST file name (without the .bst extension)

\bibliography{references} 
\rhead{
    \begin{tabular}{c}
         \includegraphics[width=3cm, height=3cm]{./logo-fa-IR.png}\\
    \end{tabular}
}
\chead{
    \setstretch{2}
    بسمه تعالی \\
    دانشکده مهندسی هوافضا \\
    \textbf{\large فرم تعریف پروژه کارشناسی ارشد}
}

\lhead{
    \ifthenelse{\isodd{\value{page}}}
        {
            \setstretch{1.5}
            \begin{tabular}{c c}
                تاریخ: & \blankField \\
                شماره: & \blankField\\
                پیوست: & \blankField                
                % \blankField & تاریخ: \\
                % \blankField & شماره:\\
                % \blankField & پیوست:
            \end{tabular}
        }
        {
            \setstretch{1.5}
            \begin{tabular}{c c}
                تاریخ: & \blankField \\
                شماره: & \blankField\\
                پیوست: & \blankField
                
                % \blankField & تاریخ: \\
                % \blankField & شماره:\\
                % \blankField & پیوست:
            \end{tabular}
        }
}

%Improvement: take advantages of the optional parameters! Currently I don't have any time to learn and implement it.
\newcommand{\Information}[9]{
{
    \setstretch{1.5}
    \small
    \bfseries
    \begin{tabular}{r r r}
         نام و نام خانوادگی : #1
         &
         شماره دانشجویی: #2
         &
         معدل: #3
         \\
         گرایش: #4
         &
         تعداد واحدهای گذرانده: #5
         &
         استاد راهنما: #6
         \\
         استاد راهنمای همکار: #7
         &
         تعداد واحد پروژه: #8
         &
         استاد ممتحن: #9
    \end{tabular}
}
}
\usepackage{array} % Add this package to use 'm' column specifier

\newcommand{\Title}[3]{
	{
		\setstretch{1.5}
		\bfseries
		\noindent
		عنوان کامل پروژه:
		
		فارسی:
        
        #1
		
		انگلیسی:


		\lr{#2}\\
		
		\noindent
		نوع پروژه:\;
		\begin{tabular}{m{10em} m{10em} m{10em}}
			کاربردی:
			\ifthenelse{#3 = 0 \or #3 = 3 \or #3 = 4}{
				$\blacksquare$
			}{
				$\square$
			}
			&
			بنیادی:
			\ifthenelse{#3 = 1 \or #3 = 3}{
				$\blacksquare$
			}{
				$\square$
			}
			&
			توسعه‌ای:
			\ifthenelse{#3 = 2 \or #3 = 3 \or #3 = 4}{
				$\blacksquare$
			}{
				$\square$
			} \\
		\end{tabular}
	}
}

% \newcounter{length}
% \newcounter{itemCount}

\newcommand{\Description}[4]{
    \setstretch{1.5}
    \noindent
    \textbf{#1:}
    \noindent
    
    #2

    \vspace{20pt}
    
    \setcounter{length}{0}
    \setcounter{itemCount}{0}
    %First obtain the length of field
    \foreach\x in#4{%
        \addtocounter{length}{1}
    }
    
    \ifthenelse{#3 = 1}{
    \textbf{کلمات کلیدی:}\;
    \foreach\x in#4{%
        \addtocounter{itemCount}{1}
        \trim{\x}
        \ifthenelse{\thelength = \theitemCount}{}{-}
    }
    \vspace{1em}
    }
}


\makeatletter
\newcommand{\ProgressTable}[1]{
{
    \noindent
    \textbf{مراحل انجام پروژه و زمان‌بندی آن:}
    \\
    
    \def\tableData{}
    
    \foreach \x\y\z in #1{
        \protected@xappto\tabledata{ \x & \y & \z\\ \noexpand\hline}
    }
    
    %Add space to each row in table:
    \renewcommand{\arraystretch}{1.5}
    \begin{center}
    \begin{tabular}{| c | m{0.8\textwidth} | c |}
        \hline
        \tabledata
    \end{tabular}
    %Add space to each row in table:
        
    \end{center}
    \renewcommand{\arraystretch}{1}
}
    
}
\makeatother

%\newcommand{\References}{
%{
%    \noindent
%    \nocite{*}
%    \textbf{مراجع:}
%    \begin{latin}
%        \printbibliography[heading=none]
%    \end{latin}
%}
%}
\newcommand{\Courses}[1]{
{
    \noindent
    \textbf{دروس مورد نیاز:}\\
    
    %Not completed yet.
    \centering
    %Add space to each row in table:
    \renewcommand{\arraystretch}{1.5}
    \begin{center}
    \begin{tabular}{| >{\centering}m{6em} | >{\centering}m{3em} | >{\centering}m{6em} | >{\centering}m{6em} | >{\centering}m{3em} | >{\centering\arraybackslash}m{6em} |}
         \hline
         \rowcolor{lightgray}
         \multicolumn{3}{|c|}{جبرانی}
         &
         \multicolumn{3}{|c|}{تخصصی
            \tiny
            (ارتباط موضوع پروژه با دروسی که دانشجو گذرانده یا باید بگذراند)
         }
         \\
         \hline
         \rowcolor{lightgray}
         گذرانده
         &
         نمره
         &
         باید بگذراند
         &
         گذرانده
         &
         نمره
         &
         باید بگذراند
         \\
         \hline
         &
         &
         &
         &
         &
         \\
         \hline
         &
         &
         &
         &
         &
         \\
         \hline
         &
         &
         &
         &
         &
         \\
         \hline
    \end{tabular}
    \end{center}
    %Add space to each row in table:
    \renewcommand{\arraystretch}{1}
}
}
\newcommand{\Signatures}{
{
    %Add space to each row in table:
    \renewcommand{\arraystretch}{1.5}
    
    \begin{center}
        \scriptsize
        \begin{tabular}{|m{15em}|m{15em}|m{15em}|}
            \hline
            استاد راهنما:
            &
            نظر گروه:
            &
            نظر کمیته تحصیلات تکمیلی دانشکده:
            \\
            تاریخ تحویل فرم به مدیر گروه:
            &
            &
            \\
            امضای استاد راهنما:
            &
            تاریخ جلسه گروه:
            &
            تاریخ جلسه کمیته:
            \\ 
            &
            امضای مدیر گروه:
            &
            امضای معاون تحصیلات تکمیلی:
            \\
            &
            &
            \\
            &
            &
            \\
            & &\\
            & &\\
            \hline
        \end{tabular}
    \end{center}
    \footnotesize
    توجه: فرم تعریف پروژه بایستی یک روز قبل از جلسه گروه توسط استاد راهنما تحویل مدیر گروه شود.
    %Add space to each row in table:
    \renewcommand{\arraystretch}{1}
}
}

%\renewcommand{\headrulewidth}{0pt}
%\renewcommand{\footrulewidth}{0pt}


\begin{document}

%First parameter: Name of student(Your name)
%Second parameter: Student Number
%Third parameter: Student's GPA
%Fourth parameter: Student's Major
%Fifth parameter: Student's passed courses credits.
%Sixth parameter: Project Supervisor.
%Seventh parameter: Project Co-Supervisor
%Eighth parameter: Project Credits.
%Ninth parameter: Project Examiner
\Information{
    علی بنی‌اسد
}{
    401209244
}{
    20
}{
    فضا
}{
    15
}{
    دکتر نوبهاری
}{}{
    6
}{
    نام استاد ممتحن
}

\noindent

%First parameter: Title of your project in Persian.
%Second parameter: Title of your project in English.
%Third parameter: Type of your project =>
%   0: Theoretical.
%   1: Practical.
%   2: Both practical & theoretical
\Title{
    هدایت بازی دیفرانسیلی با استفاده از یادگیری تقویتی مقاوم در محیط‌های پویا چندجسمی با پیشران کم}{
    Robust Reinforcement Learning Differential Game Guidance in Low-Thrust, Multi-Body Dynamical Environments
}{
    4
}

\noindent

\Description{معرفی موضوع}{
این پژوهش یک منطق هدایت مقاوم با استفاده از یادگیری تقویتی بازی‌های دیفرانسیلی 
\LTRfootnote{Differential Game Reinforcement Learning (DGRL)}
 را ارائه می‌دهد. این منطق هدایت در بستر شبکه عصبی
 \LTRfootnote{Neural Network (NN)}
  برای راهنمایی خودکار و حلقه بسته فضاپیماهایی با رانشگر کم‌پیشران در محیط‌های چند جسمی متغیر چالش‌برانگیز مانند منظومه زمین-ماه، ارائه می‌شود.
}{0}{
    {
        کلمه کلیدی نمونه
        ,
        کلمه کلیدی نمونه
        ,
        کلمه کلیدی نمونه
    }
}

\Description{اهمیت موضوع}{
مطالعه و توسعه این الگوریتم هدایت جدید با استفاده از الگوریتم‌های یادگیری تقویتی و یادگیری بازی‌های دیفرانسیلی دارای اهمیت ویژه‌ای است. این الگوریتم هدایت، به عنوان یک مدل مقاوم، قابلیت مقابله با نوسانات و متغیرهای محیطی را داراست. علاوه بر این، ساختار جدید پردازنده‌های کامپیوتر پرواز، با امکانات به‌روز و محاسبات ماتریسی بهبود یافته، اجرای بهینه‌تر الگوریتم‌های یادگیری ماتریسی را آسان کرده و توانایی اجرای موثر آنها را فراهم کرده است.

با توجه به اینکه ایستگاه زمینی در بیشتر زمان ماموریت فاصله زمانی قابل توجهی از محل اجرای دستورات تا مکان اعمال واقعی دستورات دارد، استفاده از این الگوریتم هدایت می‌تواند به عنوان یک تکنولوژی بسیار مهم در کاهش تاخیرهای ناشی از این فواصل زمانی و افزایش پایداری و دقت سیستم‌های کامپیوتر پرواز تلقی شود. توانایی این الگوریتم هدایت در انجام محاسبات پیچیده و تصمیم‌گیری‌های برنامه‌ریزی در زمان واقعی، اهمیت ویژه‌ای برای امنیت و پایداری پروازهای فضایی دارد و می‌تواند تاثیر مثبتی در اجرای موثر ماموریت‌های فضایی آینده داشته باشد.
}{0}{{}}

\Description{کاربردها}{
 کاربردهای الگوریتم هدایت در  ماموریت‌های مختلف فضایی در ادامه آورده شده است. این الگوریتم هدایت می‌تواند در ماموریت‌های فضایی مختلفی مانند ماموریت‌های ماهواره‌ای، ماموریت‌های سفر به ماه و ماموریت‌های نزدیک به زمین مورد استفاده قرار گیرد.

\begin{enumerate}
    \item \textbf{هدایت خودکار فضاپیماها:} یکی از کاربردهای اصلی این پژوهش، هدایت خودکار فضاپیماها در محیط‌های پویا و پیچیده فضایی است. این پژوهش به ایجاد یک سیستم هدایت مبتنی بر یادگیری تقویتی که توانایی انجام مانورها، تعقیب مسیر‌های مرجع بین مدارهای مختلف و بهبود اشتباهات را دارد، می‌پردازد.

    \item \textbf{حل مسائل پیچیده هدایت:} این پژوهش برای حل مسائل پیچیده هدایت فضایی که نیازمند تصمیم‌گیری در زمان واقعی و مطابق با متغیرهای محیطی می‌باشند، استفاده می‌شود. از جمله این مسائل می‌توان به تغییر مسیر در محیط‌های پیچیده، تنظیمات نیرویه پیشرانه، و اصلاح وضعیت نسبت به مدارهای مرجع اشاره کرد.

    \item \textbf{هدایت در محیط‌های نامطمئن:} از طریق استفاده از الگوریتم‌های یادگیری تقویتی، این پژوهش به هدایت در محیط‌های نامطمئن و بدون نیاز به مدل دقیق محیط فضایی امکان می‌دهد. این موضوع در مواقعی که مدل‌سازی دقیق مدل‌های دینامیکی مشکل باشد، قابل اجرا است.

    \item \textbf{استفاده در ماموریت‌های متنوع:} این پژوهش امکان استفاده در ماموریت‌های مختلف فضایی را فراهم می‌کند. از جمله ماموریت‌های مختلف می‌توان به ماموریت‌های ماهواره‌ای، تعمیر و نگهداری ماموریت فضایی، تعقیب اهداف در مدارهای مختلف، و انجام مانور‌های پیچیده اشاره کرد.
\end{enumerate}
}{0}{{}}






\Description{تعریف دقیق مسئله}{
	در پژوهش حاضر، مسأله هدایت فضاپیما برای سیستم زمین-ماه با استفاده از مدل دینامیکی سه‌جسم محدود دایره‌ای\LTRfootnote{Circular Restricted Three-Body Problem (CR3BP)} مطرح می‌شود. در این مسأله، یک فضاپیما با سیستم پیشران کم به یک مسیر مرجع در سیستم \lr{CR3BP} منتقل می‌شود. شرایط اولیه فضاپیما از مسیر مرجع انحراف‌های تصادفی دارند. هدف از این پژوهش، توسعه یک الگوریتم هدایت حلقه بسته با استفاده از اصول یادگیری تقویتی بر مبنای بازی دیفرانسیلی است که به فضاپیما اجازه می‌دهد با فرض بدترین اغتشاش به محیط مرجع بازگردد، مسیر مرجع را دنبال کرده و به مدار مقصد برسد. بازی دیفرانسیلی موجب می‌شود الگوریتم هدایت نه تنها بهینه باشد بلکه مقاوم نیز باشد.
}{0}{{}}


\Description{فرضیات مسئله}{
\begin{itemize}
	\item مدل دینامیکی مسأله، مسأله سه‌جسم محدود دایره‌ای
	سیستم زمین-ماه است. این فرضیه یک سیستم دوجسمی ساده با فضاپیما به عنوان جسم سوم بسیار کوچک را فرض می‌کند.
	\item حرکت فضاپیما تنها در صفحه اتفاق می‌افتد و تنها نیروهای گرانشی در نظر گرفته می‌شوند. سایر انحراف‌های مانند فشار تابشی خورشید نادیده گرفته می‌شوند.
	\item فضاپیما به یک سیستم پیشران کم مجهز است که شعاع تاثیر مشخصی دارد. میزان و جهت پیشران به صورت مستمر تغییر کرده و قابل تنظیم است.
	\item مسأله هدایت شامل انتقال بین مدارهای دوره‌ای (مدارهای لیاپانوف) در مدل \lr{CR3BP} زمین-ماه است.
	\item  توانایی موتور دارای محدودیت است. موتور قادر به تولید پیشران در هر جهت و مقدار نیست.
	\item وضعیت فضاپیما به طور کامل قابل مشاهده و در دسترس الگوریتم هدایت در هر مرحله زمانی است.
	\item محیط یادگیری تقویتی ویژگی‌های مارکوف را دارد، به این معنا که وضعیت فعلی تمام اطلاعات لازم برای پیش‌بینی وضعیت‌های آینده را فراهم می‌کند.
	\item آموزش در یک محیط چرخه‌ای \LTRfootnote{episodic}
	 انجام می‌شود.
	\item تابع پاداش به صورت بازی دیفرانسیلی طراحی شده  است که فضاپیما در مسیر نزدیک به مرجع با فرض بدترین اغتشاش  بماند و به مدار مقصد برسد.
	\item هدف اصلی تنها  اجرا پذیری مسأله   نیست و دیگر اهداف مانند مصرف کمتر سوخت یا زمان انتقال با استفاده از نظریه بازی در نظر گرفته می‌شوند.
	\item این رویکرد نیاز به دسترسی به منابع محاسباتی با عملکرد بالا برای آموزش یادگیری تقویتی دارد قبل از اجرای آن در رایانه پرواز فضاپیما.
\end{itemize}

	
}{0}{{}}

\Description{روش انجام کار}{
	% مرحله 1: ایجاد محیط دینامیکی
% مراحل کلی
% مراحل کلی
\begin{itemize}
	\item \textbf{مرحله 1: ایجاد محیط دینامیکی}
	\begin{itemize}
		\item پیاده‌سازی مدل دینامیکی محیط مورد نظر انجام می‌شود. این مدل باید شامل تمام جوانب مهم محیط رباشد، از جمله قوانین حرکت، ابعاد فضایی، و وضعیت اولیه.
		\item در صورت پیچیدگی مدل دینامیکی محیط، از کتابخانه‌ها و ابزارهای مهندسی نرم‌افزاری مانند NumPy و SciPy در پایتون برای پیاده‌سازی آن استفاده کرد.
		\item مدل دینامیکی محیط  با ویژگی اینکه به درستی کار کند و اطلاعات مورد نیاز برای یادگیری تقویتی (مانند وضعیت فعلی) در دسترس باشند، ایجاد می‌شود.
	\end{itemize}
	
	\item \textbf{مرحله 2: ایجاد بازیگر هدایت}
	\begin{itemize}
		\item در این مرحله، بازیگری برای انجام هدایت در محیط خود ایجاد می‌شود. این بازیگر مسئول انتخاب اعمال (اعمال کم پیشران) بر اساس وضعیت فعلی محیط است.
		\item از الگوریتم‌های یادگیری تقویتی مختلف مانند \lr{Q-Learning}، \lr{Deep Q-Networks (DQN)}، یا \lr{Proximal Policy Optimization (PPO)} برای ایجاد بازیگر هدایت استفاده می‌شود. انتخاب الگوریتم بستگی به  محیط و اهداف آموزش دارد.
	\end{itemize}
	
	\item \textbf{مرحله 3: استفاده از منابع \lr{GPU} در \lr{Google Colab}}
	\begin{itemize}
		\item \lr{Google Colab} به عنوان یک محیط آموزش آنلاین به کار می‌رود که اجازه می‌دهد کدهای پایتون در یک محیط مبتنی بر مرورگر اجرا شوند. از این مزیت این محیط می‌توان به دسترسی به منابع \lr{GPU} برای آموزش مدل‌های عمیق اشاره کرد.
		\item برای استفاده از \lr{GPU} در \lr{Google Colab}، می‌توانید از کتابخانه‌های معروفی مانند \lr{TensorFlow} یا \lr{PyTorch} که از \lr{GPU} پشتیبانی می‌کنند، استفاده می‌شود. به این ترتیب، آموزش مدل‌های یادگیری تقویتی بسیار سریع‌تر و رایگان انجام می‌شود.
	\end{itemize}
\end{itemize}

	
}{0}{{}}

\Description{پیشینه‌ی موضوع}{
	رایج‌ترین تکنیک‌های هدایت معمولاً با استفاده از ایستگاه‌های کنترل مستقر در زمین انجام می‌شوند. با این حال، حساسیت به شکست ارتباطات، تاخیرهای زمانی، محدودیت‌های انتقال داده، پیچیدگی سنسورها و هزینه‌های عملیاتی همگی از دلایلی هستند که هدایت، مسیریابی و کنترل\LTRfootnote{Guidance, Navigation and Control (GNC)}
	 را از کامپیوتر ایستگاه زمینی به کامپیوتر پرواز منتقل می‌کنند. 
در حالی که در طراحی مسیر از افزایش قدرت سخت‌افزارهای کامپیوتر بهره‌برداری می‌کنند، تعداد کمی از آن‌ها برای اجرای خودکار درون‌سفینه عملی هستند. این به دلیل محدودیت‌های منابع محاسباتی است که در فضاپیماها وجود دارد.

در فرآیند طراحی مسیر، معمولاً یک مسیر بهینه و تاریخچه فرمان کنترلی طراحی می‌شود که با معیارهای ماموریت، مصرف سوخت و زمان پرواز، مطابقت داشته باشد. این روش قبل از پرواز انجام می‌شود و می‌تواند از استراتژی‌های متعددی برای هدایت بهینه با پیشران کم استفاده کند، از جمله تکنیک‌های بهینه‌سازی جهانی \cite{vavrina2017global} و برنامه‌نویسی غیرخطی \cite{ocampo2004finite}.
}{0}{{}}
\cite{vavrina2017global} 


%\Description{پژوهش‌های داخلی}{
%}{0}{{}}
%
%
%\Description{پژوهش‌های داخل دانشگاه و دانشکده}{
%}{0}{{}}





%Table of progress:
%   Pass line number, works and their duration in a comma delimited list.
%   line number / work / duration ,
%   example: 1 / read X / 1 year ,

\ProgressTable{
    {
    1/پایداری و کنترل: برقراری پایداری و کنترل مناسب در حین ناوبری گروهی از پرندهها بسیاراهمیت دارد. هماهنگی در حرکت و جهتدهی به پرندهها بهمنظور جلوگیری از تصادف و تضاد در  حرکت ضروری است.
    \newline
    علی
    /۴ ماه,
    ۲/مرحله دوم/۳ ماه,
    ۳/مرحله سوم/۲ ماه,
    ۳/مرحله چهارم/۵ ماه,
    ۵/مرحله پنجم/۲ ماه
    }
}



%\printbibliography

%\bibliography{refs}


\noindent

% Add our references in references,bib
%\References{}
%
% -------------------------------------------------------
%  Bibliography
% -------------------------------------------------------


\begin{latin}
%	\renewcommand{\bibname}{\rl{مراجع}}
\baselineskip=.8\baselineskip
\bibliographystyle{./styles/packages/unsrtabbrv}

% Uncomment next line to include uncited references
% \nocite{*}

\bibliography{bibs/full,bibs/refs}
 
\end{latin}
\newpage


%\begin{thebibliography}{1}
%	
% -------------------------------------------------------
%  Bibliography
% -------------------------------------------------------


\begin{latin}
%	\renewcommand{\bibname}{\rl{مراجع}}
\baselineskip=.8\baselineskip
\bibliographystyle{./styles/packages/unsrtabbrv}

% Uncomment next line to include uncited references
% \nocite{*}

\bibliography{bibs/full,bibs/refs}
 
\end{latin}
\newpage

%\end{thebibliography}

\begin{latin}
	\begin{thebibliography}{1}
		
		\bibitem{vavrina2017global}
		Vavrina, Matthew A, Jacob A Englander, Sean M Phillips, and Kyle M Hughes.
		"Global, multi-objective trajectory optimization with parametric spreading."
		In AAS AIAA Astrodynamics Specialist Conference 2017, Tech. No. GSFC-E-DAA-TN45282, 2017.
		
		\bibitem{ocampo2004finite}
		Ocampo, Cesar.
		"Finite Burn Maneuver Modeling for a Generalized Spacecraft Trajectory Design and Optimization System."
		Annals of the New York Academy of Sciences 1017 (2004): 210-233.
		doi: 10.1196/annals.1311.013.
		
	\end{thebibliography}
\end{latin}


\noindent

% Change \newcommand{\Courses} in the previous lines
\Courses{}

\noindent

\Signatures{}

\end{document}
